\chapter{LITERATURE REVIEW}
Since its initial scientific investigations, hydroponics, a soilless farming method utilizing
nutrient-rich solutions has undergone substantial development. The concept that plants
can complete their life cycle with roots immersed in a balanced nutrient solution,
provided sufficient aeration is maintained, was established by Gericke’s seminal work
in 1937, which demonstrated the commercial viability of “aquaculture” \cite{gericke1937}.
This work laid the foundation for modern Controlled Environment Agriculture (CEA).\\[10pt]
Since then, hydroponic technology has evolved massively and is now widely adopted in
greenhouses, vertical farms, and indoor cultivation of plants. The limitations of conventional
soil-based agriculture are reduced as CEA integrates environmental control, nutrient
management, and automation. Savvas and Gruda \cite{savvas2018} emphasized that soilless
cultivation technologies play a critical role in modern greenhouse production by enabling
higher crop productivity, uniform growth , and reduced dependency on soil quality.\\[10pt]
One of the primary advantages of hydroponics over conventional soil-based farming is its
superior resource efficiency. Barbosa et al. \cite{barbosa2015} conducted a comparative
analysis of lettuce production using hydroponic and conventional farming methods that
showed hydroponics can reduce water consumption by 70-90\% due to nutrient
solution recirculation, significantly minimizing losses from evaporation and percolation.
Furthermore, the direct delivery of nutrients to the root zone improves fertilizer use
efficiency and reduces nutrient leaching into the environment, addressing major ecological
concerns associated with traditional agriculture \cite{savvas2018}.\\[10pt]
Another major benefit of hydroponics is the improved efficiency of nutrient utilization. In
soil-based farming, nutrients may become unavailable due to fixation, leaching, or microbial
competition. In contrast, hydroponics delivers nutrients directly to the root zone in soluble
ionic forms, allowing precise control over nutrient composition and concentration \cite{trejo2012}.
This direct nutrient delivery reduces fertilizer wastage and minimizes environmental
pollution caused by excess nutrient runoff, which is a major concern in conventional agriculture
\cite{savvas2018}.\\[10pt]
The pH and electrical conductivity (EC) of the nutrient solution are critical parameters
affecting hydroponic performance. Nutrient availability and ion solubility are strongly
influenced by pH, with most crops performing optimally within a slightly acidic range
of 5.5 to 6.5 \cite{trejo2012}. EC provides an indication of the total dissolved salts in the
solution, reflecting overall nutrient concentration. Maintaining these parameters within
narrow, crop-specific limits is essential to prevent nutrient lockout, deficiencies, or toxicities,
all of which directly impact yield and crop quality \cite{hansen1978}.\\[10pt]
EC serves as an indirect measure of the TDS in the nutrient solution and provides an
estimate of overall nutrient concentration. Maintaining EC within crop-specific limits is
essential to avoid osmotic stress, nutrient deficiencies, or toxicities. Low EC values may
result in inadequate nutrient supply, whereas excessively high EC levels can damage
plant roots, reduce water uptake, and negatively impact plant growth and yield \cite{hansen1978}.\\[10pt]
Nutrient management in hydroponics is made more difficult by the relationship
between pH and EC. The pH of a solution can be affected by changes in nutrient concentration,
and nutritional ion balance can be changed by pH alterations. Therefore, to maintain ideal
growing conditions and guarantee steady plant development, accurate and ongoing
monitoring of these parameters is required \cite{hansen1978}.\\[10pt]
Traditionally, monitoring and adjusting these parameters has relied on manual intervention,
which is time-consuming, susceptible to human error, and often inconsistent. Early
automation efforts primarily focused on timer-based nutrient dosing. However, advancements
in low-cost, reliable sensors and microcontrollers have enabled closed-loop feedback
control systems. Mendoza et al. \cite{mendoza2020} demonstrated an automated hydroponic
system using Arduino-based pH and EC sensing combined with Proportional-Integral-Derivative
(PID) control and peristaltic pumps, resulting in more stable growing conditions compared
to manual control. Similarly, Nam et al. \cite{nam1996} contrasted conventional farming
with hydroponics, aeroponics, and aquaponics, emphasizing the critical role automation plays
in increasing accuracy, lowering human error, and guaranteeing consistent root-zone
conditions. Continuous monitoring and prompt corrective intervention are made possible by
automated devices, which are especially crucial for green products like lettuce that develop
quickly.\\[10pt]
The integration of Internet of Things (IoT) represents the next evolutionary stage of
hydroponic automation. IoT-enabled projects support cloud-based data storage, real-time
remote monitoring through web or mobile interfaces, and automated warning generation.
A review by Kumar et al. \cite{kumar2021} concluded that IoT-based hydroponic frameworks
enhance traceability, provide decision-support capabilities, and facilitate the integration
of predictive control techniques, including machine learning models.\\[10pt]
IoT makes it possible to continuously monitor environmental factors, performance,
and fertilizer parameters in hydroponic applications without the need for constant personal
presence. Long-term trend analysis of pH, EC, temperature, and humidity is made possible
by data logging, which supports data-driven optimization and early anomaly detection.
IoT solutions also make it possible to provide alerts when parameters depart from permitted
ranges, which lowers the risk of crop damage and permits quick corrective action \cite{kumar2021}.\\[10pt]
Despite these advancements, several challenges remain, particularly regarding sensor reliability
and calibration drift during continuous operation. Additionally, there is a growing need for
flexible and cost-effective solutions tailored to small-scale, urban, and indoor farming
applications, as many existing methods are designed for specific crops or large-scale commercial
use.

\renewcommand{\arraystretch}{1.4}
\setlength{\tabcolsep}{3pt}
\begin{longtable}{|
p{0.65cm}|
p{1.75cm}|
p{2.2cm}|
p{1.9cm}|
p{2.05cm}|
p{2.2cm}|
p{2.3cm}|}

\caption{Comparison of Existing Systems with Designed System} \\
\hline
\textbf{S.N.} & \textbf{Reference} & \textbf{System Type} & \textbf{Control Method} & \textbf{Key Features} & \textbf{Limitations} & \textbf{Comparison with designed System} \\
\hline
\endfirsthead

\hline
\textbf{S.N.} & \textbf{Reference} & \textbf{System Type} & \textbf{Control Method} & \textbf{Key Features} & \textbf{Limitations} & \textbf{Comparison with designed System} \\
\hline
\endhead

1 & Gericke (1937) \cite{gericke1937} & Basic Hydroponics & Manual & Introduced concept of soilless farming & No automation, no monitoring system & Full automation and real-time control \\
\hline

2 & Savvas \& Gruda (2018) \cite{savvas2018} & Soilless Cultivation & Manual / Semi - automated & Improved yield and efficiency & Limited automation and control precision & Precise parameter control using sensors \\
\hline

3 & Barbosa et al. (2015) \cite{barbosa2015} & Hydroponic vs Conventional & Not control - focused & High water efficiency (70--90\% saving) & No control system implemented & Control and monitoring for efficiency \\
\hline

4 & Trejo- Téllez \& Gómez- Merino (2012) \cite{trejo2012} & Nutrient Solution Study & Manual & Defined optimal pH and nutrient ranges & No real-time monitoring or automation & Maintenance of optimal ranges \\
\hline

5 & Hansen (1978) \cite{hansen1978}& Nutrient Management & Manual & Explained EC and pH importance & No automation, prone to human error & Elimination of manual adjustment errors \\
\hline

6 & Mendoza et al. (2020) \cite{mendoza2020} & Automated Hydroponic System & PID Control & Closed- loop control using sensors and pumps & Complex tuning, risk of overshoot, higher computation & Simpler threshold- based control with stabilization delay \\
\hline

7 & Nam et al. (1996) \cite{nam1996} & Hydroponic Environment Control & Partial automation & Improved crop growth using controlled environment & Limited parameter control and no IoT integration & IoT monitoring \\
\hline

8 & Kumar et al. (2021) \cite{kumar2021} & IoT Hydroponic Systems & Varies (general review) & Cloud monitoring, remote access & Lacks optimized control strategy for low-cost systems & IoT+ efficient low-cost control method \\
\hline

9 & Designed System & Fully Automated Hydroponic with IoT & Threshold based Closed loop & Real time monitoring, automated dosing, IoT dashboard, stabilization delay & Limited to single crop, frequent calibration & -- \\
\hline

\end{longtable}